% Einheit 4 von 4 – Mit Größen rechnen im Sachzusammenhang (bank/einheiten/e4.jsonl, zusammenbau v0.1)
%% TODO Zweigzeile Teil 1: was hier gelernt wird (Satz in Schülersprache aus dem Einheitstitel) – Einheit 4
\einheitenkopf[e4]{Einheit 4 von 4 · Mit Größen rechnen im Sachzusammenhang}
\zweigzeile{ab Kl. 8 · P10}
\setcounter{aufgabe}{33}

%% TODO \verfahren{…}: Verfahrensüberschrift in Sachsprache fehlt in der Bank (2.3 b)
%% TODO Titel: Ich-kann-Satz fehlt in der Bank (Platzhalter: Kettenname) – Nr. 34
%% TODO Anweisung: ein Satz über der ersten Teilaufgabe fehlt in der Bank – Nr. 34
\begin{aufgabe}{Mit Größen rechnen im Sachzusammenhang}
%% TODO \rechenplatz{n}: Schrittzahl des Grundfalls fehlt in der Bank (nur bei fester Schreibform, 2.3 b)
\begin{teile}
\teil Von einer $1{,}2$ m langen Leiste werden $45$ cm abgesägt; gefragt ist der Rest in cm. Welche Angabe musst du vorher umrechnen? Kreuze an.\\ \kreuz{$1{,}2$ m}\\ \kreuz{$45$ cm}
\teil $1{,}4$ km $+ 650$ m – wie viel km? \leerfeld[km]
\teil Ein Modellhaus ist $0{,}08$ m hoch. Gib die Höhe in einer handlicheren Einheit an. \leerfeld \leerfeld
\teil Ein Aluminiumwürfel hat ein Volumen von $30$ cm³; Aluminium hat die Dichte $2{,}7$ g/cm³. Masse? \leerfeld[g]
\teil Ein Kupferteil wiegt $445$ g; Kupfer hat die Dichte $8{,}9$ g/cm³. Volumen? \leerfeld[cm³]
\teil Ein Wassertank fasst $2$ m³ und wiegt leer $85$ kg. Er ist zu $30\,\%$ gefüllt; ein Kubikmeter Wasser wiegt eine Tonne. Gesamtmasse? \leerfeld[kg]
\teil Eine Wand hat $18$ m²; $1$ l Farbe reicht für $6$ m². Wie viele Liter braucht man? \leerfeld[l]
\teil Zwei Wände, je $4$ m lang und $2{,}6$ m hoch, werden zweimal gestrichen; $1$ l Farbe reicht für $8$ m². Wie viele Liter? \leerfeld[l]
\teil Man braucht $1{,}3$ l Farbe. Es gibt Dosen mit $0{,}75$ l, $1$ l und $2$ l. Welche einzelne Dose reicht mindestens? \leerfeld[l]
\teil Jonas rechnet seinen Schulweg: $1{,}6$ km $+ 750$ m $= 751{,}6$ m. Nenne den Fehler und rechne richtig.
\teil Eine Messingkugel wiegt $3\,500$ g; Messing hat die Dichte $8{,}4$ g/cm³. Berechne das Volumen der Kugel. Runde auf eine Stelle nach dem Komma. \hfill (P10 2016 OS) \\ \leerfeld[cm³]
\teil Ein Vogelhaus hat zwei Dachflächen zu je $0{,}18$ m² und vier Wände zu je $0{,}12$ m². Alles wird zweimal gestrichen; $1$ l Farbe reicht für $8$ m². Es gibt Dosen mit $125$ ml, $250$ ml und $750$ ml. Lea meint, die kleinste Dose reicht, Max will die mittlere kaufen. Wer hat recht? Begründe. \hfill (P10 2015 OS)
\teil Ein Busunternehmen hat auf $120\,000$ km Mehrkosten von $744$ € für Reifen. In der Zeitung steht: „Jeder Kilometer wird dadurch etwa $6$ ct teurer.“ Die Rechnung dazu: $744$ € $= 744\,000$ ct; $744\,000 : 120\,000 = 6{,}2$ ct. Erkläre den Fehler und korrigiere die Behauptung. \hfill (P10 2014 OS)
\end{teile}
\end{aufgabe}

%% TODO \verfahren{…}: Verfahrensüberschrift in Sachsprache fehlt in der Bank (2.3 b)
%% TODO Titel: Ich-kann-Satz fehlt in der Bank (Platzhalter: Kettenname) – Nr. 35
%% TODO Anweisung: ein Satz über der ersten Teilaufgabe fehlt in der Bank – Nr. 35
\begin{aufgabe}{Kosten und Einnahmen, Sek II}
%% TODO \rechenplatz{n}: Schrittzahl des Grundfalls fehlt in der Bank (nur bei fester Schreibform, 2.3 b)
\begin{teile}
\teil Die Fläche ist in cm² gegeben, der Druckpreis in Euro pro m². Darf man direkt malnehmen? Kreuze an.\\ \kreuz{ja}\\ \kreuz{nein}
\teil Ein Kiosk verkauft $23$ Eis der Sorte A und $17$ der Sorte B, beide zu $2$ €. Einnahme? \leerfeld[€]
\teil Ein Stand verkauft Stofftaschen: weiße $22$ und schwarze $17$ zu je $6$ €, bunte $15$, $9$ und $12$ zu je $4$ €. Einnahme? \leerfeld[€]
\teil Ein Plakat hat $7\,500$ cm²; der Druck kostet $30$ € pro m². Kosten? \leerfeld[€]
\teil Aus einer $2{,}4$ m langen Latte werden zwei Leisten gesägt, $5$ und $3$ Längeneinheiten lang; eine Längeneinheit entspricht $25$ cm. Wie viel Prozent der Latte sind Verschnitt? \hfill (FHR 2022) \\ \leerfeld \%
\end{teile}
\end{aufgabe}

%% TODO Titel: Ich-kann-Satz fehlt in der Bank (Platzhalter: Kettenname) – Nr. 36
%% TODO Anweisung: ein Satz über der ersten Teilaufgabe fehlt in der Bank – Nr. 36
\begin{aufgabe}{Dichte aus Masse und Volumen (Vorrat)}
\begin{teile}
\teil Ein Stein wiegt $540$ g und hat ein Volumen von $200$ cm³. Dichte? \leerfeld g/cm³
\end{teile}
\end{aufgabe}

%% TODO Titel: Ich-kann-Satz fehlt in der Bank (Platzhalter: Kettenname) – Nr. 37
%% TODO Anweisung: ein Satz über der ersten Teilaufgabe fehlt in der Bank – Nr. 37
\begin{aufgabe}{Euro und Cent in einer Rechnung}
\begin{teile}
\teil Drei Brezeln zu je $85$ ct und ein Kakao zu $2{,}40$ € – zusammen? \leerfeld[€]
\end{teile}
\end{aufgabe}

%% TODO Titel: Ich-kann-Satz fehlt in der Bank (Platzhalter: Kettenname) – Nr. 38
%% TODO Anweisung: ein Satz über der ersten Teilaufgabe fehlt in der Bank – Nr. 38
\begin{aufgabe}{Mit Größen rechnen im Sachzusammenhang – Fehler finden, Begründen, Anwendung}
\begin{teile}
\teil Tim schreibt: $7{,}25$ € $= 7\,250$ ct. Wo ist der Fehler?
\teil Beim Sportfest werden $45$ Getränke zu $3$ € und $25$ zu $1$ € verkauft. Lina rechnet: Durchschnittspreis $2$ €, $70 \cdot 2 = 140$ €. Wo ist der Fehler?
\teil Ein Ergebnis lautet $0{,}007$ km. Begründe, warum man es besser in Metern angibt.
\teil Begründe, warum man bei $40$ Karten zu $8$ € und $5$ Karten zu $4$ € nicht mit dem Durchschnittspreis $6$ € rechnen darf.
\teil Ein Fliesenleger belegt einen Badboden mit $6{,}3$ m². Eine Fliese ist $30$ cm mal $30$ cm groß. Wie viele Fliesen braucht er mindestens (ohne Verschnitt)? \leerfeld Fliesen
\end{teile}
\end{aufgabe}
